\section{Standard Library}\label{sec:standard-library}

This chapter documents the standard library built-in functions and additional operator overloads. These features provide native support for generating random values, mathematical rounding, and specific string/array operations.

% ─────────────────────────────────────────────
\subsection{Random Generation}\label{sec:random-generation}

QLang includes globally available built-in functions for generating random values. These functions do not support keyword arguments.

\begin{itemize}
    \item \texttt{seed(value)} — Initializes the internal pseudorandom number generator with the specified numeric seed. If called with no arguments, it returns the current seed.
    \item \texttt{random()} — Returns a pseudorandom number in the range $[0.0, 1.0)$.
\end{itemize}

% ─────────────────────────────────────────────
\subsection{Mathematical Rounding}\label{sec:math-rounding}

The rounding functions take a numeric value and an optional \texttt{decimals} argument. If \texttt{decimals} is omitted or is less than or equal to $0$, the function removes any fractional part (returning a whole number). Otherwise, it returns a number rounded to the specified fractional precision.

\begin{itemize}
    \item \texttt{cut(val, decimals=0)} — Truncates the fractional part of \texttt{val} towards zero.
    \item \texttt{round(val, decimals=0)} — Rounds \texttt{val} to the nearest number. Halfway values are rounded away from zero.
    \item \texttt{floor(val, decimals=0)} — Returns the largest whole number less than or equal to \texttt{val}.
    \item \texttt{ceil(val, decimals=0)} — Returns the smallest whole number greater than or equal to \texttt{val}.
\end{itemize}

% Auto-generated by docs/tools/highlight.py - do not edit
\begin{codebox}
(*@\textcolor{codebox_comment}{// Set random seed}@*)

(*@\textcolor{black}{seed}@*)(*@\textcolor{codebox_operator}{(}@*)(*@\textcolor{codebox_number}{12345}@*)(*@\textcolor{codebox_operator}{)}@*)(*@\textcolor{codebox_operator}{;}@*)



(*@\textcolor{codebox_comment}{// Get a random float [0, 1)}@*)

(*@\textcolor{codebox_type}{\textbf{num}}@*) (*@\textcolor{black}{r}@*) (*@\textcolor{codebox_operator}{=}@*) (*@\textcolor{black}{random}@*)(*@\textcolor{codebox_operator}{(}@*)(*@\textcolor{codebox_operator}{)}@*)(*@\textcolor{codebox_operator}{;}@*)



(*@\textcolor{codebox_comment}{// Math rounding and truncating functions}@*)

(*@\textcolor{codebox_type}{\textbf{num}}@*) (*@\textcolor{black}{val}@*) (*@\textcolor{codebox_operator}{=}@*) (*@\textcolor{codebox_number}{12.3456}@*)(*@\textcolor{codebox_operator}{;}@*)



(*@\textcolor{codebox_type}{\textbf{num}}@*) (*@\textcolor{black}{c}@*) (*@\textcolor{codebox_operator}{=}@*) (*@\textcolor{black}{cut}@*)(*@\textcolor{codebox_operator}{(}@*)(*@\textcolor{black}{val}@*)(*@\textcolor{codebox_operator}{,}@*) (*@\textcolor{codebox_number}{2}@*)(*@\textcolor{codebox_operator}{)}@*)(*@\textcolor{codebox_operator}{;}@*)      (*@\textcolor{codebox_comment}{// 12.34}@*)

(*@\textcolor{codebox_type}{\textbf{num}}@*) (*@\textcolor{black}{f}@*) (*@\textcolor{codebox_operator}{=}@*) (*@\textcolor{black}{floor}@*)(*@\textcolor{codebox_operator}{(}@*)(*@\textcolor{black}{val}@*)(*@\textcolor{codebox_operator}{,}@*) (*@\textcolor{codebox_number}{1}@*)(*@\textcolor{codebox_operator}{)}@*)(*@\textcolor{codebox_operator}{;}@*)    (*@\textcolor{codebox_comment}{// 12.3}@*)

(*@\textcolor{codebox_type}{\textbf{num}}@*) (*@\textcolor{black}{cl}@*) (*@\textcolor{codebox_operator}{=}@*) (*@\textcolor{black}{ceil}@*)(*@\textcolor{codebox_operator}{(}@*)(*@\textcolor{black}{val}@*)(*@\textcolor{codebox_operator}{,}@*) (*@\textcolor{codebox_number}{0}@*)(*@\textcolor{codebox_operator}{)}@*)(*@\textcolor{codebox_operator}{;}@*)    (*@\textcolor{codebox_comment}{// 13}@*)

(*@\textcolor{codebox_type}{\textbf{num}}@*) (*@\textcolor{black}{rnd}@*) (*@\textcolor{codebox_operator}{=}@*) (*@\textcolor{black}{round}@*)(*@\textcolor{codebox_operator}{(}@*)(*@\textcolor{black}{val}@*)(*@\textcolor{codebox_operator}{,}@*) (*@\textcolor{codebox_number}{3}@*)(*@\textcolor{codebox_operator}{)}@*)(*@\textcolor{codebox_operator}{;}@*)  (*@\textcolor{codebox_comment}{// 12.346}@*)


\end{codebox}


% ─────────────────────────────────────────────
\subsection{Advanced Operator Overloading}\label{sec:adv-operator-overloading}

QLang provides implicit operator overloads for \texttt{text} variables and arrays. These operators perform data transformations, such as string manipulation and array concatenation.

\begin{center}
\begin{tabular}{|p{3cm}|p{9cm}|}
\hline
\textbf{Expression} & \textbf{Behavior} \\
\hline
\texttt{text - text} & 
    Removes all occurrences of the characters present in the right-hand string from the left-hand string. \\
\hline
\texttt{text * num} & 
    Repeats the text a specified number of times. \\
\hline
\texttt{text / text} & 
    Splits the left-hand string by the right-hand separator string, returning an array of parts. \\
\hline
\texttt{array + array} & 
    Concatenates two arrays into a single array. \\
\hline
\texttt{array + element} & 
    Appends a single element to the end of the array. \\
\hline
\end{tabular}
\end{center}

% Auto-generated by docs/tools/highlight.py - do not edit
\begin{codebox}
(*@\textcolor{codebox_type}{\textbf{text}}@*) (*@\textcolor{black}{word}@*) (*@\textcolor{codebox_operator}{=}@*) (*@\textcolor{codebox_string}{\string"hello world\string"}@*)(*@\textcolor{codebox_operator}{;}@*)
(*@\textcolor{codebox_type}{\textbf{text}}@*) (*@\textcolor{black}{clean}@*) (*@\textcolor{codebox_operator}{=}@*) (*@\textcolor{black}{word}@*) (*@\textcolor{codebox_operator}{-}@*) (*@\textcolor{codebox_string}{\string"o\string"}@*)(*@\textcolor{codebox_operator}{;}@*)       (*@\textcolor{codebox_comment}{// \string"hell wrld\string"}@*)

(*@\textcolor{codebox_type}{\textbf{text}}@*) (*@\textcolor{black}{repeated}@*) (*@\textcolor{codebox_operator}{=}@*) (*@\textcolor{codebox_string}{\string"ab\string"}@*) (*@\textcolor{codebox_operator}{*}@*) (*@\textcolor{codebox_number}{3}@*)(*@\textcolor{codebox_operator}{;}@*)      (*@\textcolor{codebox_comment}{// \string"ababab\string"}@*)

(*@\textcolor{codebox_type}{\textbf{text}}@*) (*@\textcolor{black}{path}@*) (*@\textcolor{codebox_operator}{=}@*) (*@\textcolor{codebox_string}{\string"user/local/bin\string"}@*)(*@\textcolor{codebox_operator}{;}@*)
(*@\textcolor{codebox_type}{\textbf{text}}@*) (*@\textcolor{black}{parts}@*)(*@\textcolor{codebox_operator}{[}@*)?(*@\textcolor{codebox_operator}{]}@*) (*@\textcolor{codebox_operator}{=}@*) (*@\textcolor{black}{path}@*) (*@\textcolor{codebox_operator}{/}@*) (*@\textcolor{codebox_string}{\string"/\string"}@*)(*@\textcolor{codebox_operator}{;}@*)    (*@\textcolor{codebox_comment}{// [\string"user\string", \string"local\string", \string"bin\string"]}@*)

(*@\textcolor{codebox_type}{\textbf{num}}@*) (*@\textcolor{black}{a}@*)(*@\textcolor{codebox_operator}{[}@*)?(*@\textcolor{codebox_operator}{]}@*) (*@\textcolor{codebox_operator}{=}@*) (*@\textcolor{codebox_operator}{[}@*)(*@\textcolor{codebox_number}{1}@*)(*@\textcolor{codebox_operator}{,}@*) (*@\textcolor{codebox_number}{2}@*)(*@\textcolor{codebox_operator}{]}@*)(*@\textcolor{codebox_operator}{;}@*)
(*@\textcolor{codebox_type}{\textbf{num}}@*) (*@\textcolor{black}{b}@*)(*@\textcolor{codebox_operator}{[}@*)?(*@\textcolor{codebox_operator}{]}@*) (*@\textcolor{codebox_operator}{=}@*) (*@\textcolor{codebox_operator}{[}@*)(*@\textcolor{codebox_number}{3}@*)(*@\textcolor{codebox_operator}{,}@*) (*@\textcolor{codebox_number}{4}@*)(*@\textcolor{codebox_operator}{]}@*)(*@\textcolor{codebox_operator}{;}@*)
(*@\textcolor{codebox_type}{\textbf{num}}@*) (*@\textcolor{black}{combined}@*)(*@\textcolor{codebox_operator}{[}@*)?(*@\textcolor{codebox_operator}{]}@*) (*@\textcolor{codebox_operator}{=}@*) (*@\textcolor{black}{a}@*) (*@\textcolor{codebox_operator}{+}@*) (*@\textcolor{black}{b}@*)(*@\textcolor{codebox_operator}{;}@*)       (*@\textcolor{codebox_comment}{// [1, 2, 3, 4]}@*)
(*@\textcolor{codebox_type}{\textbf{num}}@*) (*@\textcolor{black}{appended}@*)(*@\textcolor{codebox_operator}{[}@*)?(*@\textcolor{codebox_operator}{]}@*) (*@\textcolor{codebox_operator}{=}@*) (*@\textcolor{black}{a}@*) (*@\textcolor{codebox_operator}{+}@*) (*@\textcolor{codebox_number}{5}@*)(*@\textcolor{codebox_operator}{;}@*)       (*@\textcolor{codebox_comment}{// [1, 2, 5]}@*)

\end{codebox}

